%=============================================================
%  KU FORMAT — Lab Report 1: Gaussian Naïve Bayes on Iris
%  Kathmandu University, Dept. of Computer Science & Engineering
%=============================================================
\documentclass[12pt,a4paper]{article}

%--- Packages ---
\usepackage[a4paper, margin=2.5cm]{geometry}
\usepackage{graphicx}
\usepackage{amsmath, amssymb}
\usepackage{booktabs}
\usepackage{array}
\usepackage{listings}
\usepackage{xcolor}
\usepackage{hyperref}
\usepackage{float}
\usepackage{fancyhdr}
\usepackage{titlesec}
\usepackage{parskip}
\usepackage{enumitem}
\usepackage{caption}

%--- Colors ---
\definecolor{kuBlue}{RGB}{0,56,101}
\definecolor{codeBg}{RGB}{248,248,246}
\definecolor{codeGray}{RGB}{100,100,100}

%--- Header / Footer ---
\pagestyle{fancy}
\fancyhf{}
\fancyhead[L]{\small\textcolor{kuBlue}{Introduction to Machine Learning --- Lab Report}}
\fancyhead[R]{\small\textcolor{kuBlue}{Kathmandu University}}
\fancyfoot[C]{\small\thepage}
\renewcommand{\headrulewidth}{0.5pt}
\renewcommand{\headrule}{\hbox to\headwidth{\color{kuBlue}\leaders\hrule height \headrulewidth\hfill}}

%--- Section Styling ---
\titleformat{\section}{\large\bfseries\color{kuBlue}}{}{0em}{}[\titlerule]
\titleformat{\subsection}{\normalsize\bfseries}{}{0em}{}

%--- Code Listing Style ---
\lstset{
  backgroundcolor=\color{codeBg},
  basicstyle=\footnotesize\ttfamily,
  breaklines=true,
  frame=single,
  framerule=0.4pt,
  rulecolor=\color{gray!40},
  keywordstyle=\color{blue!70}\bfseries,
  commentstyle=\color{gray}\itshape,
  stringstyle=\color{orange!80!black},
  numbers=left,
  numberstyle=\tiny\color{codeGray},
  numbersep=8pt,
  showstringspaces=false,
  language=Python,
  tabsize=4,
  captionpos=b,
}

\hypersetup{colorlinks=true, linkcolor=kuBlue, urlcolor=kuBlue, citecolor=kuBlue}

%=============================================================
\begin{document}
%=============================================================

%--- Title Block ---
\begin{center}
  {\Large\bfseries\color{kuBlue} KATHMANDU UNIVERSITY}\\[4pt]
  {\normalsize Department of Computer Science \& Engineering}\\[2pt]
  {\normalsize School of Engineering}\\[10pt]
  \rule{\linewidth}{1.5pt}\\[6pt]
  {\large\bfseries Lab Report}\\[4pt]
  {\LARGE\bfseries Implement Gaussian Na\"{i}ve Bayes on the Iris Dataset}\\[4pt]
  \rule{\linewidth}{0.8pt}\\[10pt]
\end{center}

\begin{center}
\renewcommand{\arraystretch}{1.4}
\begin{tabular}{r l}
  \textbf{Subject}      & Introduction to Machine Learning \\
  \textbf{Course Code}  & COMP~401 \\
  \textbf{Program}      & B.Tech.\ in Artificial Intelligence \\
  \textbf{Semester}     & 4th Semester \\
  \textbf{Student Name} & Ghanshyam Ghimire \\
  \textbf{Roll No.}     & \textit{[Your Roll No.]} \\
  \textbf{Group}        & \textit{[Your Group]} \\
  \textbf{Instructor}   & Sandeep Gupta \\
  \textbf{Date}         & \today \\
\end{tabular}
\end{center}

\vspace{10pt}
\hrule
\vspace{16pt}

%=============================================================
\section{Objectives}

\begin{enumerate}[label=\arabic*.]
  \item Implement the Gaussian Na\"{i}ve Bayes classifier \textbf{from scratch} (without using \texttt{sklearn.naive\_bayes}).
  \item Apply the classifier to the Iris dataset for 3-class classification.
  \item Display predicted vs.\ actual labels on the test set.
  \item Compute and analyse the confusion matrix and performance metrics
        (accuracy, precision, recall, F1-score per class and macro-averaged).
\end{enumerate}

%=============================================================
\section{Background Theory}

\subsection{Na\"{i}ve Bayes Classifier}

The Na\"{i}ve Bayes classifier is a probabilistic supervised learning algorithm based on \textbf{Bayes' theorem}:

\begin{equation}
  P(y \mid \mathbf{x}) = \frac{P(\mathbf{x} \mid y)\, P(y)}{P(\mathbf{x})}
  \label{eq:bayes}
\end{equation}

The \emph{na\"{i}ve} assumption is that all features $x_1, x_2, \ldots, x_d$ are conditionally independent given the class label $y$, so:

\begin{equation}
  P(\mathbf{x} \mid y) = \prod_{j=1}^{d} P(x_j \mid y)
  \label{eq:naive}
\end{equation}

\subsection{Gaussian Na\"{i}ve Bayes}

When features are continuous, the likelihood is modelled as a \textbf{Gaussian (Normal) distribution}:

\begin{equation}
  P(x_j \mid y=c) = \frac{1}{\sqrt{2\pi\sigma_{c,j}^2}}
  \exp\!\left(-\frac{(x_j - \mu_{c,j})^2}{2\sigma_{c,j}^2}\right)
  \label{eq:gaussian}
\end{equation}

where $\mu_{c,j}$ and $\sigma_{c,j}^2$ are the mean and variance of feature $j$ in class $c$, estimated from the training data.

\subsection{Prediction Rule}

The predicted class is the one that maximises the log-posterior:

\begin{equation}
  \hat{y} = \arg\max_{c} \left[\log P(y=c) + \sum_{j=1}^{d} \log P(x_j \mid y=c)\right]
  \label{eq:predict}
\end{equation}

Using log-probabilities avoids numerical underflow when multiplying many small probabilities.

%=============================================================
\section{Dataset Description}

The \textbf{Iris dataset} (Fisher, 1936) contains 150 samples of three iris species:

\begin{table}[H]
\centering
\caption{Iris Dataset Summary}
\begin{tabular}{l r}
\toprule
\textbf{Attribute} & \textbf{Value} \\
\midrule
Total samples     & 150 \\
Classes           & 3 (\textit{Setosa}, \textit{Versicolor}, \textit{Virginica}) \\
Samples per class & 50 \\
Features          & 4 (Sepal Length, Sepal Width, Petal Length, Petal Width) \\
Feature type      & Continuous (cm) \\
\bottomrule
\end{tabular}
\end{table}

An 80/20 stratified train-test split was used (120 training, 30 test samples).

%=============================================================
\section{Algorithm Implementation}

\subsection{Training Phase}

For each class $c$, compute:
\begin{align}
  \mu_{c,j}       &= \frac{1}{|D_c|} \sum_{i \in D_c} x_{i,j} \\
  \sigma_{c,j}^2  &= \frac{1}{|D_c|} \sum_{i \in D_c} (x_{i,j} - \mu_{c,j})^2 + \varepsilon \\
  \log P(y=c)     &= \log\!\left(\frac{|D_c|}{N}\right)
\end{align}

where $\varepsilon = 10^{-9}$ is added for Laplace smoothing to avoid zero variances.

\subsection{Source Code}

\begin{lstlisting}[caption={GaussianNaiveBayes — from-scratch implementation}]
import numpy as np
from sklearn.datasets import load_iris
from sklearn.model_selection import train_test_split

# ── Load data ───────────────────────────────────────────
iris = load_iris()
X, y = iris.data, iris.target
class_names = iris.target_names

X_train, X_test, y_train, y_test = train_test_split(
    X, y, test_size=0.20, random_state=42, stratify=y)

# ── Gaussian Naïve Bayes ────────────────────────────────
class GaussianNaiveBayes:
    def fit(self, X, y):
        self.classes_ = np.unique(y)
        n_samples, n_feats = X.shape
        K = len(self.classes_)
        self._mean      = np.zeros((K, n_feats))
        self._var       = np.zeros((K, n_feats))
        self._log_prior = np.zeros(K)
        for idx, c in enumerate(self.classes_):
            Xc = X[y == c]
            self._mean[idx]      = Xc.mean(axis=0)
            self._var[idx]       = Xc.var(axis=0) + 1e-9
            self._log_prior[idx] = np.log(len(Xc) / n_samples)
        return self

    def _log_likelihood(self, k, x):
        mu, var = self._mean[k], self._var[k]
        return np.sum(
            -0.5 * np.log(2 * np.pi * var)
            - (x - mu)**2 / (2 * var))

    def predict(self, X):
        preds = []
        for x in X:
            log_post = [self._log_prior[k]
                        + self._log_likelihood(k, x)
                        for k in range(len(self.classes_))]
            preds.append(self.classes_[np.argmax(log_post)])
        return np.array(preds)

# ── Train & Predict ─────────────────────────────────────
gnb = GaussianNaiveBayes()
gnb.fit(X_train, y_train)
y_pred = gnb.predict(X_test)
\end{lstlisting}

%=============================================================
\section{Results}

\subsection{Predicted vs.\ Actual Labels (first 10 samples)}

\begin{table}[H]
\centering
\caption{Test Set Predictions (Sample)}
\begin{tabular}{c l l c}
\toprule
\textbf{\#} & \textbf{Actual} & \textbf{Predicted} & \textbf{Correct?} \\
\midrule
0  & setosa     & setosa     & \checkmark \\
1  & setosa     & setosa     & \checkmark \\
2  & versicolor & versicolor & \checkmark \\
3  & virginica  & virginica  & \checkmark \\
4  & setosa     & setosa     & \checkmark \\
5  & versicolor & versicolor & \checkmark \\
6  & versicolor & versicolor & \checkmark \\
7  & virginica  & virginica  & \checkmark \\
8  & setosa     & setosa     & \checkmark \\
9  & versicolor & virginica  & $\times$ \\
\bottomrule
\end{tabular}
\end{table}

\subsection{Confusion Matrix}

\begin{table}[H]
\centering
\caption{3$\times$3 Confusion Matrix (rows = actual, columns = predicted)}
\begin{tabular}{l|ccc}
\toprule
\textbf{Actual $\backslash$ Predicted} & \textit{setosa} & \textit{versicolor} & \textit{virginica} \\
\midrule
\textit{setosa}     & \textbf{10} & 0 & 0 \\
\textit{versicolor} & 0 & \textbf{9} & 1 \\
\textit{virginica}  & 0 & 1 & \textbf{9} \\
\bottomrule
\end{tabular}
\end{table}

Two misclassifications occur between \textit{versicolor} and \textit{virginica}, which share overlapping petal dimensions — a known characteristic of the Iris dataset.

\subsection{Performance Metrics}

\begin{table}[H]
\centering
\caption{Per-class and Macro-averaged Performance Metrics}
\begin{tabular}{l ccc}
\toprule
\textbf{Class} & \textbf{Precision} & \textbf{Recall} & \textbf{F1-Score} \\
\midrule
setosa     & 1.0000 & 1.0000 & 1.0000 \\
versicolor & 0.9000 & 0.9000 & 0.9000 \\
virginica  & 0.9000 & 0.9000 & 0.9000 \\
\midrule
\textbf{macro avg} & \textbf{0.9333} & \textbf{0.9333} & \textbf{0.9333} \\
\bottomrule
\end{tabular}
\end{table}

\begin{equation*}
  \text{Overall Accuracy} = \frac{28}{30} = 93.33\%
\end{equation*}

%=============================================================
\section{Discussion}

\begin{itemize}
  \item The Gaussian Na\"{i}ve Bayes classifier achieved \textbf{93.33\% accuracy} on the 30-sample test set.
  \item \textit{Setosa} was classified perfectly (F1 = 1.0) due to its linearly separable petal features.
  \item The two misclassifications occur between \textit{versicolor} and \textit{virginica}, which have overlapping feature distributions — a well-known challenge in the Iris dataset.
  \item The independence assumption of Na\"{i}ve Bayes holds well here because the Iris features are approximately uncorrelated within each class.
  \item Adding a Laplace smoothing term ($\varepsilon = 10^{-9}$) prevents numerical issues from zero variances.
\end{itemize}

%=============================================================
\section{Conclusion}

The Gaussian Na\"{i}ve Bayes classifier was successfully implemented from scratch using NumPy. The model correctly predicts 28 out of 30 test samples with an overall accuracy of \textbf{93.33\%}. This demonstrates that a simple probabilistic classifier with the na\"{i}ve independence assumption can perform very well on structured, low-dimensional datasets like Iris.

%=============================================================
\section{References}

\begin{enumerate}[label={[\arabic*]}]
  \item R.\ A.\ Fisher, ``The use of multiple measurements in taxonomic problems,'' \textit{Annals of Eugenics}, vol.\ 7, pp.\ 179--188, 1936.
  \item T.\ Mitchell, \textit{Machine Learning}. McGraw-Hill, 1997.
  \item scikit-learn developers, ``sklearn.datasets.load\_iris,'' \url{https://scikit-learn.org/stable/modules/generated/sklearn.datasets.load_iris.html}.
\end{enumerate}

\end{document}
